% Score de recherche de chacune des 37 questions, réparti selon le périmètre.
% Le point de la figure est la position du seuil : il sépare proprement la
% bande basse, et pas les deux bandes hautes — c'est délibéré.
%
% Les trois bandes forment une série catégorielle (bleu, ambre, prune) et non
% une échelle de statut : une question hors périmètre n'est pas un échec, elle
% relève d'un autre domaine. Le vert et le rouge, réservés au succès et à
% l'échec ailleurs dans le rapport, sont donc écartés ici.
% Score de recherche par question, réparti en trois bandes
% Généré par rapport/build_data.py — NE PAS ÉDITER.
\newcommand{\ScoresTravail}{(0.5955,2.055) (0.6039,1.890) (0.6053,1.945) (0.6131,2.000) (0.6328,2.055) (0.6542,2.110) (0.6611,1.890) (0.6716,1.945) (0.6743,2.000) (0.6793,2.055) (0.6797,2.110) (0.6920,1.890) (0.6924,1.945) (0.6938,2.000) (0.6946,2.055) (0.6958,2.110) (0.6960,1.890) (0.7050,1.945) (0.7076,2.000) (0.7092,2.055) (0.7124,2.110) (0.7135,1.890) (0.7153,1.945) (0.7160,2.000) (0.7164,2.055) (0.7173,2.110) (0.7180,1.890) (0.7214,1.945) (0.7263,2.000) (0.7278,2.055) (0.7335,2.110) (0.7373,1.890) (0.7434,1.945) (0.7487,2.000) (0.7492,2.055) (0.7550,2.110) (0.7558,1.890) (0.7640,1.945) (0.7658,2.000) (0.7670,2.055) (0.7674,2.110) (0.7707,1.890) (0.7805,1.945) (0.7826,2.000) (0.7840,2.055) (0.7881,2.110) (0.7935,1.890)}
\newcommand{\ScoresJuridique}{(0.4453,1.055) (0.5147,0.890) (0.5150,0.945) (0.5561,1.000) (0.6027,1.110)}
\newcommand{\ScoresHorsDroit}{(0.3071,-0.110) (0.3357,-0.055) (0.4319,0.000) (0.4461,0.110)}

\begin{tikzpicture}
\begin{axis}[
  bandesY,
  height=5.8cm,
  xmin=0.26, xmax=0.85,
  ymin=-0.75, ymax=3.05,
  xtick={0.3,0.4,0.5,0.6,0.7,0.8},
  xticklabel style={/pgf/number format/fixed,
                    /pgf/number format/precision=1},
  xlabel={score de recherche \textcolor{slate}{(similarité cosinus maximale)}},
  ytick={0,1,2},
  yticklabels={},
  legend style={at={(0.5,1.1)}, anchor=south},
]
% Seuil : filet vertical plein hauteur, étiquette posée en haut pour ne
% croiser ni les graduations ni les libellés de zone.
\draw[anrtred, line width=0.8pt, dash pattern=on 3pt off 2pt]
  (axis cs:0.42,-0.75) -- (axis cs:0.42,3.05);
\node[font=\sffamily\scriptsize\bfseries, text=anrtred, anchor=south east]
  at (axis cs:0.415,3.05) {seuil \SeuilAbstention~};

\addplot[only marks, mark=*, mark size=1.9pt,
         draw=inptblue, fill=inptblue, fill opacity=0.7]
  coordinates {\ScoresTravail};
\addplot[only marks, mark=*, mark size=1.9pt,
         draw=amberdark, fill=amberdark, fill opacity=0.8]
  coordinates {\ScoresJuridique};
\addplot[only marks, mark=*, mark size=1.9pt,
         draw=plum, fill=plum, fill opacity=0.8]
  coordinates {\ScoresHorsDroit};
\legend{droit du travail, autre domaine juridique, hors du droit}

% Étiquettes de bande, à gauche du nuage : elles remplacent les graduations
% de l'axe vertical, qui ne porte aucune quantité.
\node[font=\sffamily\scriptsize, text=inptblue!70!black, anchor=east]
  at (axis cs:0.272,2) {\NbInScope{} questions};
\node[font=\sffamily\scriptsize, text=amberdark!70!black, anchor=east]
  at (axis cs:0.272,1) {3 questions};
\node[font=\sffamily\scriptsize, text=plum!70!black, anchor=east]
  at (axis cs:0.272,0) {2 questions};

% Les deux zones que le seuil sépare réellement, sous le nuage.
\node[font=\sffamily\scriptsize\itshape, text=anrtred, anchor=north east]
  at (axis cs:0.412,-0.42) {rejeté avant le modèle};
\node[font=\sffamily\scriptsize\itshape, text=slate, anchor=north west]
  at (axis cs:0.428,-0.42) {transmis au modèle, qui peut encore s'abstenir};
\end{axis}
\end{tikzpicture}
